\section{Experimental Setup}
\label{sec:setup}

Four studies, presented here in the order their results appear: the staged
\emph{contention study} (\S\ref{sec:setup-staged}), the \emph{topology and
capacity study} of agentic RAG (\S\ref{sec:setup-topology}), two
\emph{instrument-fidelity} experiments that calibrate the harness itself
(\S\ref{sec:setup-fidelity}), and the MLPerf \emph{load-pattern reduction}
(\S\ref{sec:setup-reduction}). This section specifies designs, metrics, and
--- for measurements still in flight at submission --- hypotheses with their
falsification outcomes; no empirical claim in this section is a result.

\subsection{Devices Under Test}
\label{sec:setup-duts}

The \textbf{Apple M2~Pro} is a \SI{16}{\giga\byte} unified-memory laptop SoC
whose CPU, GPU, and Neural Engine (ANE) share one physical memory at
${\approx}\SI{200}{\giga\byte\per\second}$; LLMs run as 4-bit MLX
weights~\cite{mlx}, the ANE via Core~ML. The \textbf{NVIDIA GB10}
(DGX~Spark~\cite{gb10}) couples a Grace CPU and a Blackwell GPU over a
coherent NVLink-C2C interconnect to \SI{128}{\giga\byte} of LPDDR5X at
\SI{273}{\giga\byte\per\second} --- one physical pool, but with an
asymmetric access path: the GPU reaches memory across the interconnect.
Models run as BF16 PyTorch weights. We never
treat the pair as a controlled axis: capacity differs $8\times$, bandwidth
${\approx}1.4\times$, and device class, runtime, and precision co-vary ---
including arithmetic intensity: at batch~1, 4-bit weights carry
${\approx}4\times$ more compute per weight byte than BF16, so both decode
configurations are bandwidth-bound in absolute terms but to different
degrees. Cross-device statements therefore
concern the direction of an effect on each platform separately; absolute
comparisons stay within a device, and each platform is one physical machine
($n{=}1$ per platform). Environments are pinned in versioned specifications;
the drivers hard-abort on mismatched environments; every run records its
git commit and environment fingerprint.

\subsection{Metrics, Statistics, and Pre-Specified Configuration}
\label{sec:setup-metrics}

\paragraph{Latency anchor.} Reported per-query latency is \emph{response
time}: completion minus the query's \emph{intended} arrival time from the
load schedule. Anchoring at intended arrival (rather than at dequeue or even
enqueue) means queueing delay --- including interference-induced delay ---
is inside the reported number, and a saturated cell shows up as unbounded
growth rather than as flattering service times; this is the same discipline
the arrival log enforces on the offered-load
side~\cite{tene2015coordinated}. Throughput is completed queries over
makespan. Stage-level metrics: \textbf{stage concurrency} (sweep-line count
of temporally overlapping stage intervals --- overlap, not CPU parallelism:
pure-Python stages serialize on the GIL, native kernels overlap) and
\textbf{per-device occupancy} (union of run intervals per engine ---
residency, not utilization).

\paragraph{Memory traffic.} On the M2~Pro we sample the memory controller's
per-requestor read/write byte counters via a custom IOReport-based sampler
(no privileges required). Because the requestor-to-engine mapping is not
complete on this SoC --- under a pure-GPU load a fraction of traffic reports
under fabric agents --- we calibrate the mapping once per engine with
single-engine loads and report the residual unattributed fraction alongside
every attribution. Per-engine is not per-pipeline: when two pipelines drive
the same engine (Step A's foreground and background both use the GPU), the
counters report their sum, and we attribute by measuring the background's
schedule-determined traffic in isolation and subtracting, with the
subtraction's error carried into the reported interval. On GB10 the DRAM profiling counter is unavailable on
current drivers; we use the memory-controller busy fraction plus
model-based traffic estimates (weights plus KV-cache bytes per operation
$\times$ operation rate; activations excluded and stated), and mark every
GB10 traffic figure \emph{proxy-backed}. Thermal state is
monitored per run (package power and temperature); runs crossing a
pre-specified throttle threshold are excluded and counted. On the
\SI{16}{\giga\byte} device we additionally record memory pressure and swap
activity per cell and require quiescence --- distinguishing bandwidth
contention from capacity pressure, which would otherwise masquerade as a
nonlinear dose--response (\S\ref{sec:setup-staged}).

\paragraph{Quality.} Answer quality is exact match (EM) and token-level
F1~\cite{rajpurkar2016squad} with SQuAD normalization, over 120 questions
per arm in dedicated serial runs. All answer-producing arms share one
literal short-answer instruction (defined once, imported by every prompt
formatter), validated before collection. We do not claim ``matched''
quality from overlapping intervals: quality comparisons are \emph{paired on questions} --- every arm answers
the same 120 questions, so the EM difference is a paired proportion
(McNemar-style interval), roughly halving the width an unpaired comparison
would need --- and are reported as \emph{bounded differences}: a speed
result is called quality-preserving only when the paired interval excludes
differences beyond a pre-specified non-inferiority margin of 10 points
(chosen for decidability at $n{=}120$ and stated as such); otherwise it is
a trade-off. Containment (substring) appears only as a secondary
signal beside mean answer length; token-Jaccard parity checks that arms
answer the same task (agreement, not correctness).

\paragraph{Statistics.} The replication unit is the run. Cheap cells run
$R{=}10$, expensive cells $R{=}5$, fixed distinct seeds; we report pooled
medians with hierarchical bootstrap CIs (runs resampled, then queries;
$10^4$ resamples), print raw per-run values beside every interval, and
treat $R{=}5$ intervals as approximate --- key contrasts also report the
paired across-run difference where the design pairs (interleaving, shared
run slots). Wilson intervals cover rate metrics (EM, answered); F1, a
bounded mean, uses the hierarchical bootstrap. p95 appears only where the
pooled post-warm-up sample reaches 500 queries; p99 never. Queries hitting
a cell timeout are reported as failures with their count, never silently
dropped from latency distributions. Generation is greedy everywhere,
verified by repeat-run identical-output checks.

\paragraph{Pre-specified configuration.} Every load-generation knob derives
from a committed rule with pilot-measured inputs, locked into the
configurations before the corresponding collection. We say
\emph{pre-specified}, not pre-registered: rules, derivations, and configs
are committed and timestamped in the artifact's version history, not
registered with a third party. The rules: open-loop rates are $\rho =
0.6\times$ the pilot-measured capacity ($1/\mathrm{median}$ serial service
time) of the slowest \emph{piloted core arm} of each comparison, one rate
per task per device; the contention foreground uses its own stated rule,
$0.4\times$ capacity, so the baseline sits mid-utilization with headroom
for degradation to be visible; queue depths equal the query count, so the
entry queue cannot throttle arrivals --- verified per run from the
intended-vs-realized arrival log (zero blocked submissions; realized rate
within 5\% of the drawn schedule). Warm-up drops come from a rolling-median
flatness detector with a $2\times$ margin \emph{where the detector
converges}; where a pilot is too short or heavy-tailed to converge, we fall
back to a stated default (two queries for LLM pipelines), mark the cell,
and report a per-experiment warm-up sensitivity check rather than trusting
either path blindly. Pilot capacity estimates are single short serial runs
and are treated accordingly: every open-loop cell carries a post-hoc
stability check (no trend in queue depth across the measurement window),
and cells failing it are reported at reduced scope rather than as latency
results. Table~\ref{tab:knobs-topology} shows one experiment's knob table; the
artifact carries all of them with derivations, pilot traces, and
verification outcomes.

\subsection{The Contention Study (Steps A--D)}
\label{sec:setup-staged}

The capability claim under test: collocated pipelines are first-class and
attributable. The design descends from a production-shaped symptom to a
candidate mechanism in four stages, each changing \emph{exactly one}
configuration element of its predecessor; the paper shows each diff.

\textbf{Step A (system view).} The foreground is plain RAG serving
(retrieve top-3 from a vector index; generate with a 4B model) at a fixed
Poisson rate. Beside it run $B \in \{0,1,2\}$ independent \emph{indexing}
pipelines (chunk $\rightarrow$ GPU embed $\rightarrow$ vector-store
insert), each its own OS process with its own disjoint corpus shard and its
own target collection. The foreground never reads the background's
collections, so index growth cannot alter foreground behavior; the planned
check is byte-identical foreground answers across $B$ under greedy
decoding, making hardware the only coupling. Metrics: foreground
p95/median response time and throughput vs.\ $B$ ($\geq 500$ pooled
post-warm-up foreground queries per cell), background throughput,
per-process attribution, memory-traffic counters, and the
capacity-pressure telemetry above.

\textbf{Step B (isolate intensity).} $B{=}1$; \emph{only the background
loadgen block} changes --- saturating $\rightarrow$ fixed-interval at
$\{25,50,75,100\}\%$ of the indexer's pilot-measured maximum rate.

\textbf{Step C (purify the co-runner).} \emph{Only the background stage
list} changes: the indexer is replaced by single-resource co-runners at
matched intensity ladders --- a calibrated CPU streaming stage
(STREAM-triad~\cite{mccalpin1995stream} over a working set far beyond the
last-level cache; its bytes per operation are deterministic, so its
offered traffic is known without any counter), a GPU
CLIP~\cite{radford2021clip} encode loop, and, on the M2~Pro, the same
encoder on the ANE. GPU and ANE co-runner traffic comes from the
memory-controller counters on the M2~Pro and from model-based estimates on
GB10; the CPU streaming co-runner, whose traffic is analytic, anchors
both --- decoupling the dose axis from the instrument that measures the
response.

\textbf{Step D (purify the foreground).} \emph{Only the foreground}
changes: RAG serving becomes bare LLM decode, and analysis splits
time-to-first-token (prefill: processing the prompt) from per-token decode
(generating output). Decode streams the weights plus the growing KV cache
every token and is bandwidth-dominated at batch~1; prefill is
compute-dominated at these prompt lengths on the 4-bit
device~\cite{patel2024splitwise,agrawal2024sarathi}, while on GB10 (BF16)
short prompts sit nearer the compute/bandwidth ridge, so the control's
contrast is sharp on the M2~Pro and attenuated on GB10 --- we state its
strength per platform. Both phases are measured as \emph{service} time
(stage start to end): the control isolates where contention lands inside a
query, so queueing delay is deliberately outside it, unlike the
response-time metrics of Steps A--B.

\emph{Hypotheses and outcomes.} H1: foreground degradation grows with
co-runner memory traffic at matched intensity, largely independently of
the generating engine --- tested cleanly on the co-runners that do
\emph{not} share the foreground's engine (CPU stream; ANE on the M2~Pro).
The GPU co-runner shares the foreground's engine, and on GB10 runs in a
separate process where CUDA contexts time-slice the device (MPS is not
enabled), so its arm confounds device time with bandwidth; we report it as
a labeled composite, not as an H1 test. Slope compatibility is judged
against a pre-specified equivalence band (slope ratios within
$[2/3, 3/2]$), not by interval overlap. H2: degradation concentrates in decode; prefill
moves substantially less. Falsification branches are first-class outcomes:
engine-\emph{dependent} slopes at matched traffic would localize contention
in engine-specific paths (fabric, schedulers) rather than DRAM; a flat
dose--response would bound the platforms' co-runner tolerance. Slopes are fitted only
below the saturation knee; behavior near the roofline is reported
descriptively, since memory-controller queueing is not linear there.
\emph{Further threats:} occupancy is residency; attribution is
per-engine-calibrated on the M2~Pro and proxy-backed on GB10; $B$ is
bounded by device memory, and the ceiling is reported, not extrapolated.

\subsection{The Topology and Capacity Study}
\label{sec:setup-topology}

One self-reflective RAG task --- grade, generate, check, rewrite-on-failure,
in the style popularized by Self-RAG~\cite{asai2024selfrag} and corrective
RAG~\cite{yan2024crag}, implemented as explicit pipeline stages rather than
reflection tokens --- and four topologies that swap by configuration
alone. \textbf{Monolith-9B}: one Qwen3.5-9B~\cite{qwen35}
grades retrieved context, answers, and self-checks in one structured pass,
with a bounded rewrite cycle. \textbf{Decomposed}: three 4B specialists
(grader, generator, hallucination-checker), so cheap calls for one query
can overlap expensive generation for another. \textbf{Monolith-4B}: one 4B
issuing the three calls sequentially --- the size-matched control.
\textbf{Decomposed-Shared}: the decomposed graph over \emph{one} resident
4B behind a mutex, separating graph structure from physical residency.
Because stages share a default device stream, model-level GPU parallelism
is unavailable by construction (\S\ref{sec:framework-execution}); any
Decomposed-over-Shared gap therefore isolates what extra \emph{residency}
buys under scheduling-level overlap (unserialized CPU phases, no mutex
convoying), not concurrent kernel execution, and we interpret it as such.
All execution backends here are batch-1 engines; production servers
overlap queries \emph{within} one instance via continuous
batching~\cite{yu2022orca,kwon2023vllm}, so this study's conclusions are
about graph-structured batch-1 pipelines and are not serving guidance ---
a served-engine backend exists in the framework and extending the
comparison to it is future work.
The specialists share base weights differentiated by prompt; resident
memory is reported as unique plus shared mapped pages, since memory-mapped
weights can share physical pages across instances. All rungs of the
capacity sweep are dense models.
We run two difficulties: factoid (rag-mini-wikipedia, top-3 over 3{,}200 passages) and multi-hop (HotpotQA~\cite{yang2018hotpotqa}, top-5 over a 20{,}000-passage corpus of gold plus distractor passages via ChromaDB~\cite{chroma}). Retrieval is fixed within a difficulty and contrasts are read within a difficulty; the size control runs on \emph{both} tasks, since a multi-hop
quality gap is exactly what a size difference alone would produce. A
top-$k{=}10$ sensitivity column guards the multi-hop result against
retrieval-depth artifacts.

Cells: \emph{timing} (40 queries, pipelined and serial, $R{=}5$; all arms
of a task share one open-loop rate) and \emph{quality} (120 questions,
serial) are separate runs, each sized for its statistic. Latency
comparisons come from the open-loop cells at the shared rate;
\emph{service-capacity} comparisons --- the throughput headline --- come
from the closed-loop serial cells, where completion rate is the arm's own
service rate rather than the offered load. \textbf{Capacity sweep.} The
monolith arm read along the model-size axis --- $\{0.8, 2, 4, 9, 27\}$B,
each rung one \texttt{model}-field edit --- reports resident memory and
quality per size with each device's ceiling marked: the
\SI{15.7}{\giga\byte} 27B rung exceeds the M2~Pro's budget by
construction; GB10 runs it with headroom.

\subsection{Instrument Fidelity}
\label{sec:setup-fidelity}

\textbf{CAL-1 (harness overhead).} Pass-through pipelines isolate the harness:
a depth sweep ($1$--$100$ no-op stages, closed loop) for dispatch cost and
its scaling; a payload sweep (\SI{0}{}--\SI{10}{\mebi\byte}; reference
passing vs.\ per-hop deep copy) for the data path. M2~Pro only: with no
kernel in the loop this measures framework and interpreter, and one quiet
machine bounds it; CAL-1's absolute numbers are platform-scoped, and the GB10
fidelity claim rests on CAL-2, which runs on both devices. No-op stages never
release the GIL, making the per-stage figure a worst case; the copy arm
lower-bounds a serializing design.

\textbf{CAL-2 (modularity overhead).} The same workload two ways: a
hand-written PyTorch script vs.\ the equivalent \sysname{} pipeline ---
EfficientNetV2-S~\cite{tan2021efficientnetv2} transfer-learning
(fine-tuning a frozen backbone) on Imagenette~\cite{imagenette}, batch 8,
identical data loading, matched \texttt{num\_workers}=0, a device
synchronize closing each step. Interleaved arms, paired statistic,
$R{=}10$, 1{,}100 steps per run, 200-step warm-up (detector-confirmed).
Two instrumentation layers are measured separately: core dispatch (tracing
off) and the tracing layer (asynchronous span export). A wrapper can only
add work, so a negative estimate is treated as a confound detector, not a
result --- a stance that has already caught two real confounds
(non-monotonic cross-thread clocks; un-interleaved schedules with
background prefetch).

\subsection{The Load-Pattern Reduction}
\label{sec:setup-reduction}

MLPerf Inference's four scenarios are, at the load-generation level, four
arrival processes: one-in-flight (SingleStream), Poisson at a target rate
(Server), fixed-interval batched queries (MultiStream as specified through
MLPerf Inference v1.1; since v2.0 the scenario issues back-to-back
8-sample queries, a closed-loop pattern our SingleStream scheduler covers
at batch~8 --- we reproduce the fixed-interval form and say so), and
all-at-once (Offline). We instantiate each as a \sysname{} scheduler configuration over
one vision pipeline (ResNet-50), $R{=}10$ per scenario per device, and
verify each from its arrival log: realized rate within 5\% and exponential
inter-arrivals by quantile--quantile inspection (Server); interval
regularity (MultiStream); single-outstanding-query (SingleStream);
enqueue-before-start (Offline). The claim is scoped: \emph{load patterns},
not certified-harness semantics --- latency bounds, early stopping, and
accuracy modes are MLPerf's compliance machinery, not its arrival
processes, and are out of scope. We likewise substitute an
auto-downloadable ImageNet subset for the license-gated original: the
reduction concerns load, not accuracy. The harness footprint is
re-examined \emph{in this regime} rather than assumed from CAL-2: at
millisecond service times, \SI{60}{\micro\second}-class stage costs are
percent-level, and we report them beside the scenario metrics.

\begin{table}[t]\centering\small
\caption{Knob table for the topology study (M2~Pro / GB10). Every value
derives from a committed rule with pilot inputs; verification outcomes and
the remaining experiments' tables ship in the artifact. Rates in
queries/s.}
\label{tab:knobs-topology}
\resizebox{\columnwidth}{!}{\begin{tabular}{lllll}
\toprule
Knob & M2~Pro & GB10 & Rule & Verified by \\
\midrule
rate (factoid) & 0.117 & 0.154 & $0.6\times$ slowest piloted arm & arrival log \\
rate (multi-hop) & 0.0217 & 0.046 & $0.6\times$ slowest piloted arm & arrival log \\
timing queries & \multicolumn{2}{c}{40 $\times$ $R{=}5$} & 30--50 per cell & stability check \\
quality questions & \multicolumn{2}{c}{120, serial} & Wilson width $\leq 9$\,pts & --- \\
queue depth & \multicolumn{2}{c}{$=$ queries} & never blocks & arrival log \\
warm-up $k$ & \multicolumn{2}{c}{2 (fallback; detector n/c)} & flatness $\times 2$, else default & sensitivity check \\
top-$k$ & \multicolumn{2}{c}{3 / 5 (factoid / multi-hop)} & task design ($+k{=}10$ sens.) & --- \\
\bottomrule
\end{tabular}}
\end{table}
